%% 第26回編集委員会（物理第4巻）：事務局。第3巻の ph3-front.tex と同じ部品・同じ形・同じページ数（各2ページ）
%% ph4-front.tex  前付の見開き2つ（物理シリーズ第4巻）
%%   地図は物理シリーズ全体のもの（第1〜3巻と同じ表）。この巻の章の札を前に出し、運動量の行に「電磁波の放射圧」（電-15 ⑤）、
%%   エネルギーの行に「ジュール熱」（電-4 ⑤）、電荷の行に「連続の式」（電-4 ⑥・電-14 ④）を足した。
%%   相似の表は、第1〜3巻の前付と波-1 ⑩・電-9 ⑨の表（ph:電-9:souji）の5行と同じ文・同じ記号（Y-3：力 F ↔ 電圧 e、速度 v ↔ 電流 i、
%%   m ↔ L、c ↔ R、1/k ↔ C）。電W4 の提案（docs/作業メモ/物理第4巻_組どうしの確かめ.md）で「磁気回路」の列を足した。
%%   列の中身は電-7 ⑤の表（ph:電-7:tab-souji）と同じ：押すもの 起磁力 NI、流れるもの 磁束 Φ、流れにくさ 磁気抵抗 𝓡_m。
%%   磁気回路の相似は「押すもの・流れるもの・流れにくさ」の3つだけで、磁気抵抗は熱にせずエネルギーをためる（電-7 ⑤の「せいかくに言うと」）
%%   ことを、性格の囲みの2つめに書いた。出会う章に電-7（磁気回路）・電-9（回路）を足した。
%%   相似の図は第3巻と同じ（ばね・おもり・ダンパーと RLC の直列回路。電-9 ⑨・波-1 ⑩・6-7⑦）。
%%   この本の使い方の「式の読み方」の見本は、電力の式 P＝VI（電-4 ⑤）。数は 100 V・1200 W のドライヤーの 12 A・72 kJ（電-4 ⑤・例題1、熱-1）。
%%   道しるべの見本は、電-4 の道しるべ（p4-04.tex の \bkmichishirube）から選んだ札。
%%   やくそくに、第26回の安全の決まり（計画書 0 の8・手引き 2章の12：コンセント・高電圧を使う実験はしない）を4つめとして足した。
%% 第27回 確認の直し（修4-1、E1-19・E1-20・E1-23）：地図の運動量の行に第1〜3巻と同じ「運動量保存」をもどし、エネルギーの行の
%%   「入る・出るもの」から、法則・章にない「波」をけずった（第1・2巻の行と同じ形）。相似の表の見出しは第1〜3巻と同じ文のまま
%%   「力学」と「（ばね・おもり）」の2行に折り、colsep を第1〜3巻と同じ 3pt にした。磁気回路の列は「―」を「（ない）」に、
%%   磁気抵抗に「（失わない）」を添えた（性格の囲みの「熱を出さない」と合わせる）。
%%   p.4〜5  保存則と相似の地図  \phConservationMapPages
%%   p.6〜7  この本の使い方      \phHowToPages
%% 章のページ（\phchappage{電-2}）は ph4-pages.tex（gen_phys.py が作る）から読む。
\makeatletter
% 前付の見出し（数学の「この本の路線図」「この本の使い方」と同じ形）
\newcommand\ph@fronthead[2]{%
  \noindent{\sffamily\bfseries\fontsize{24}{30}\selectfont\color{bkDeep}#1}\hfill{\sffamily\fontsize{8.5}{11}\selectfont\color{bkGray}#2}\par\vspace{1mm}
  {\color{bkVol}\hrule height 1.4pt}\vspace{3.5mm}}
% 章の札（ページつき）。この巻（第4巻）の章はページも出す。ほかの巻の章は札だけ
\newcommand\ph@cp[1]{\ph@routestop{#1}}
\newcommand\ph@cps[1]{\foreach \c [count=\k] in {#1} {\ifnum\k>1 \hspace{0.25em}\fi\ph@cp{\c}}}

%% ============================================================
%% p.4〜5  保存則と相似の地図
%% ============================================================
\newcommand\phConservationMapPages{%
\clearpage\setcounter{page}{4}%
\global\bk@decofalse\bkfolio{true}\bksetrunning{}{}{}%
{\def\R##1##2{##1}\pdfbookmark[0]{保存則と相似の地図}{ph-front-map}\typeout{BKPAGE front 保存則と相似の地図 \thepage}}%
\ph@fronthead{保存則と相似の地図}{全5巻の背骨}
{\fontsize{10.5}{19}\selectfont
物理のたくさんの法則の多くは、たった1つの形の式から生まれます。それが下の「バランス式」で、おこづかい帳のお金の出し入れと同じ形です。\par}
\vspace{1mm}
\begin{bkfigure}
\begin{tikzpicture}[x=1mm, y=1mm]
  % 見えない境界（検査体積）
  \draw[bkTheme, line width=1.1pt, dash pattern=on 3pt off 2pt, rounded corners=3mm, fill=bkTheme!5] (38,0) rectangle (102,27);
  \node[anchor=north west, text=bkDeep, font=\sffamily\bfseries\footnotesize] at (40,26) {見えない境界で囲んだ所};
  % たまる量（水位）
  \fill[bkBlue!22] (40,2) rectangle (100,13);
  \draw[bkBlue, line width=1pt] (40,13) -- (100,13);
  \node[text=bkBlue!70!black, font=\sffamily\bfseries\small] at (70,7.5) {たまる量（増える・減る）};
  % 生じる量
  \node[draw=phFlowR, line width=0.9pt, fill=white, star, star points=8, star point ratio=1.6, minimum size=6mm, inner sep=0pt] (g) at (55,18) {};
  \node[anchor=west, text=phFlowR!80!black, font=\sffamily\bfseries\small] at (60,18) {生じる量};
  % 入る・出る
  \draw[-{Stealth[length=3.6mm, width=3mm]}, line width=3.2pt, bkTheme] (6,15) -- (37,15);
  \node[anchor=south, text=bkDeep, font=\sffamily\bfseries\small] at (20,17) {入る量};
  \draw[-{Stealth[length=3.6mm, width=3mm]}, line width=3.2pt, bkGray] (103,15) -- (134,15);
  \node[anchor=south, text=bkInk, font=\sffamily\bfseries\small] at (119,17) {出る量};
\end{tikzpicture}
\end{bkfigure}
\vspace{1mm}
\begin{center}
\tikz\node[fill=bkTheme, text=white, rounded corners=2mm, inner xsep=5mm, inner ysep=2.6mm, font=\sffamily\bfseries\large]
  {入る量 \,$-$\, 出る量 \,$+$\, 生じる量 \,$=$\, たまる量};
\end{center}
\vspace{1mm}
{\fontsize{10.5}{19}\selectfont
「数える量」をかえると、力学・流体・熱・電気・原子炉の法則になります。生じる量が0の量は「保存する量」です。札の横の数は、この巻のページです。\par}
\vspace{1.5mm}
{\footnotesize\noindent
\begin{tblr}{colspec={Q[l,m,wd=19mm] Q[l,m,wd=25mm] Q[l,m,wd=17mm] X[l,m] Q[l,m,wd=36mm]}, width=\linewidth,
  row{1}={bg=bkVol!14, font=\sffamily\bfseries\footnotesize}, hlines={0.4pt,bkLine}, rowsep=0.5pt, colsep=2.5pt,
  column{1}={font=\sffamily\bfseries}}
数える量 & 入る・出るもの & 生じる量 & 法則・式 & 章 \\
質量 & 質量流量 & 0 & 連続の式 & \ph@cps{力-13,熱-4} \\
運動量 & 力・運動量の流れ & 0 & 運動方程式、運動量保存、運動量の法則、放射圧 & {\ph@cps{電-15,力-2}\\ \ph@cps{力-3,力-13}} \\
角運動量 & トルク & 0 & 回転の運動方程式、ポンプのオイラーの式 & \ph@cps{力-6,力-13} \\
エネルギー & 熱・仕事・エンタルピー・電磁場の流れ & 0 & ポインティングベクトル、ジュール熱、力学的エネルギー保存、第一法則、ベルヌーイの式 & {\ph@cps{電-4,電-14}\\ \ph@cps{力-4,熱-2}} \\
電荷 & 電流 & 0 & 電荷保存（連続の式）、キルヒホッフの電流則 & \ph@cps{電-4,電-14} \\
水分 & 水蒸気の流れ・結露・加湿 & 0 & 水分のバランス（空気線図） & \ph@cps{熱-9} \\
エントロピー & 熱÷温度 & 0以上 & 第二法則 & \ph@cps{熱-6,熱-15} \\
中性子 & 中性子の流れ & 生まれる数－吸収される数 & 中性子のバランス、増倍率 & \ph@cps{原-15} \\
\end{tblr}\par}
\clearpage
%% ---------------- p.5 相似 ----------------
\ph@fronthead{相似（にた形の式）の地図}{分野が違っても、式は同じ形}
{\fontsize{10.5}{19}\selectfont
ばねにつけたおもりのゆれと、コイル・コンデンサー・抵抗の回路の電気のゆれは、同じ形の式で書けます（一方を解けば、もう一方も解ける）。分野が違っても式が同じ形になることを、\textbf{相似}とよびます。\par}
\vspace{1mm}
\begin{bkfigure}
\begin{tikzpicture}[x=1mm, y=1mm]
  \path[use as bounding box] (0,-17.5) rectangle (139,32);
  % 左：ばね・ダンパー・おもり
  \fill[bkLine!60] (0,4) rectangle (3,26);
  \draw[bkInk, line width=0.8pt] (3,4) -- (3,26);
  \foreach \y in {6,10,...,26} \draw[bkInk, line width=0.5pt] (0,\y) -- (3,\y-3);
  \draw[bkInk, line width=0.9pt, decorate, decoration={zigzag, amplitude=1.6mm, segment length=3mm}] (3,20) -- (31,20);
  \node[anchor=south, font=\sffamily\footnotesize] at (17,22.5) {ばね $k$};
  \draw[bkInk, line width=0.9pt] (3,9) -- (12,9) (22,5) -- (12,5) -- (12,13) -- (22,13) (17,9) -- (31,9);
  \draw[bkInk, line width=1.6pt] (17,6.3) -- (17,11.7);
  \node[anchor=north, font=\sffamily\footnotesize] at (17,4) {ダンパー $c$};
  \filldraw[fill=bkTheme!20, draw=bkTheme, line width=1pt, rounded corners=1mm] (31,4) rectangle (47,26);
  \node[font=\sffamily\bfseries\small] at (39,15) {$m$};
  \draw[-{Stealth[length=3mm]}, line width=1.4pt, phFlowR] (47,15) -- (60,15);
  \node[anchor=south, font=\sffamily\footnotesize] at (54,16.5) {力 $F$};
  \node[anchor=north, align=center, font=\small] at (31,-1) {$m\ddot x+c\dot x+kx=F$\\[1mm]{\footnotesize $x$：おもりのずれ\quad $\omega_{\mathrm n}=\sqrt{k/m}$}};
  % まん中の等号
  \node[font=\sffamily\bfseries\Large, text=bkDeep] at (70,15) {＝};
  \node[font=\sffamily\bfseries\footnotesize, text=bkDeep] at (70,8) {同じ形};
  % 右：RLC の直列回路
  \begin{scope}[shift={(79,0)}]
    \draw[bkInk, line width=0.9pt] (6,4) -- (6,11) (6,19) -- (6,26) -- (16,26);
    \draw[bkInk, line width=0.9pt] (16,23.5) rectangle (28,28.5);
    \node[anchor=south, font=\sffamily\footnotesize] at (22,29) {抵抗 $R$};
    \draw[bkInk, line width=0.9pt] (28,26) -- (34,26);
    \draw[bkInk, line width=0.9pt] (34,26) \foreach \i in {0,...,3} { arc (180:0:1.6) } -- (52,26);
    \node[anchor=south, font=\sffamily\footnotesize] at (40.4,28.5) {コイル $L$};
    \draw[bkInk, line width=0.9pt] (52,26) -- (52,17.5) (49,17.5) -- (55,17.5) (49,15.5) -- (55,15.5) (52,15.5) -- (52,4) -- (6,4);
    \node[anchor=west, font=\sffamily\footnotesize] at (56,16.5) {$C$};
    \draw[bkInk, line width=0.9pt] (6,15) circle (4);
    \node[font=\footnotesize] at (6,15) {$e$};
    \draw[-{Stealth[length=2.6mm]}, line width=1.2pt, phFlowR] (20,6.5) -- (36,6.5);
    \node[anchor=north, font=\sffamily\footnotesize] at (28,5.2) {電流 $i$};
  \end{scope}
  \node[anchor=north, align=center, font=\small] at (111,-1) {$L\ddot q+R\dot q+\dfrac{q}{C}=e$\\[1mm]{\footnotesize $q$：たまった電気\quad $\omega_{\mathrm n}=1/\sqrt{LC}$}};
\end{tikzpicture}
\end{bkfigure}
\vspace{1mm}
{\footnotesize\noindent
\begin{tblr}{colspec={Q[l,m,wd=17mm] X[l,m] X[l,m] X[l,m] X[l,m] X[l,m]}, width=\linewidth,
  row{1}={bg=bkVol!14, font=\sffamily\bfseries\footnotesize}, hlines={0.4pt,bkLine}, rowsep=0.6pt, colsep=3pt,
  column{1}={font=\sffamily\bfseries}}
役わり & {力学\\（ばね・おもり）} & {電気\\（回路）} & {磁気回路\\（鉄心）} & {熱\\（部屋・お湯）} & {流体\\（U字管）} \\
流れるもの & 速度 $v$ & 電流 $i$ & 磁束 $\Phi$ & 熱流 $\dot Q$ & 体積流量 $Q$ \\
押すもの & 力 $F$ & 電圧 $e$ & 起磁力 $NI$ & 温度差 $\Delta T$ & 圧力差 $\Delta p$ \\
ためるもの & ばねのやわらかさ $1/k$ & 静電容量 $C$ & （ない） & 熱容量 $C_{\mathrm{th}}$ & 液面の広さ $S/(\rho g)$（液面1つ） \\
慣性 & 質量 $m$ & インダクタンス $L$ & （ない） & （ない） & 管の中の液 $\rho l/S$ \\
失うもの & ダンパー $c$ & 抵抗 $R$ & （失わない）流れにくさ：磁気抵抗 $\mathcal R_{\mathrm m}$ & 熱抵抗 $R_{\mathrm{th}}$ & 管のこすれ \\
\end{tblr}\par}
\vspace{1mm}
{\footnotesize\noindent\textbf{出会う章}\quad \ph@cp{電-7}（磁気回路）\quad \ph@cp{電-9}（回路）\quad \ph@cp{波-1}（ゆれ）\quad 解き方（6-7・6-13）\par}
\begin{seikaku}
\begin{itemize}[leftmargin=1.2\zw, itemsep=0pt, topsep=0pt, label={\color{bkNavy}●}]
\item 力学の列は「力と電圧」が組です。力×速度と電圧×電流が、どちらも仕事率だからです（電-9 ⑨）。だから $\sqrt{k/m}$ は $1/\sqrt{LC}$ と読みかえられます。
\item 磁気回路の列は、押すもの・流れるもの・流れにくさの3つだけの相似です（電-7 ⑤）。磁気抵抗は熱を出さず、鉄心の磁場にエネルギーをためます。
\item 熱には慣性がないので、部屋の温度はゆれずに近づきます。
\end{itemize}
\end{seikaku}
\clearpage\global\bk@decotrue}

%% ============================================================
%% p.6〜7  この本の使い方（物理の本）
%% ============================================================
\newcommand\ph@howstep[7]{% x, y, アイコン, 色, 名前, 説明, ノード名
  \node[draw=#4, line width=1pt, fill=white, rounded corners=2.5mm,
    minimum width=30.5mm, minimum height=33mm, anchor=north] (#7) at (#1,#2) {};
  \node[anchor=north] (#7i) at ([yshift=-2.2mm]#7.north) {\bkicon[1.7em]{#3}{#4}};
  \node[anchor=north, font=\sffamily\bfseries\small, text=#4, inner sep=0pt, align=center,
    minimum height=8.5mm, text width=29.5mm] (#7l) at ([yshift=-1mm]#7i.south) {#5};
  \node[anchor=north, font=\sffamily\footnotesize, text width=28mm, align=center, text=bkInk, inner sep=0pt]
    at ([yshift=-1.0mm]#7l.south) {#6};}
% 新しい囲みの見本の札  \ph@newbox{アイコン}{色}{名前}{説明}
\newcommand\ph@newbox[4]{%
  \begin{tikzpicture}[x=1mm, y=1mm]
    \draw[#2, line width=0.9pt, fill=white, rounded corners=2mm] (0,0) rectangle (67,14);
    \node[anchor=west, inner sep=0pt] (i) at (2.5,7) {\bkicon[1.5em]{#1}{#2}};
    \node[anchor=north west, inner sep=0pt, text=#2!75!black, font=\sffamily\bfseries\small] at (11.5,12) {#3};
    \node[anchor=north west, inner sep=0pt, text width=54mm, align=left, text=bkInk, font=\sffamily\fontsize{8.5}{12}\selectfont] at (11.5,6.6) {#4};
  \end{tikzpicture}}
\newcommand\phHowToPages{%
\clearpage\setcounter{page}{6}%
\global\bk@decofalse\bkfolio{true}\bksetrunning{}{}{}%
{\pdfbookmark[0]{この本の使い方}{ph-front-howto}\typeout{BKPAGE front この本の使い方 \thepage}}%
\ph@fronthead{この本の使い方}{}
{\fontsize{10.5}{19}\selectfont
この本は、身の周りの「なぜ？」に物理で答える本です。できごとを、何かを無視した「モデル」にして式を立て、答えを数で出して確かめます。式は数学シリーズの道具で書きます。\par}

\bksubsection{1}{1つの章は、いつも同じ順番}
\begin{bkfigure}
\begin{tikzpicture}[x=1mm, y=1mm]
  \ph@howstep{0}{0}{hatena}{bkYellow!85!black}{はてな？}{身近なふしぎ\\から始めよう}{s1}
  \ph@howstep{35}{0}{yattemiyou}{bkGreen}{やってみよう}{実験して、\\はかってみよう}{s2}
  \ph@howstep{70}{0}{phmodel}{phModel}{モデル化の札}{何を無視するかを\\決めよう}{s3}
  \ph@howstep{105}{0}{phflow}{bkTheme}{一般の式から、\\よく使う式へ}{条件で、よく使う\\式を導く}{s4}
  \ph@howstep{0}{-38}{reidai}{bkBlue}{例題}{7つの段で、\\ていねいに解く}{s5}
  \ph@howstep{35}{-38}{tsumazuki}{bkRed}{まちがえやすい\\ところ}{理由を知ろう}{s6}
  \ph@howstep{70}{-38}{phgear}{bkPurple}{工学ではこう使う}{本物の機械の\\数で確かめよう}{s7}
  \ph@howstep{105}{-38}{check}{bkTheme}{章末問題}{基本とチャレンジ\\の2段だよ}{s8}
  \foreach \a/\b in {s1/s2,s2/s3,s3/s4,s5/s6,s6/s7,s7/s8}
    {\draw[bk arrow, bkGray, shorten <=1mm, shorten >=1mm] (\a.east) -- (\b.west);}
  \draw[bk arrow, bkGray, rounded corners=2mm, shorten <=1.5mm, shorten >=1.5mm] (s4.south) -- ++(0,-2.5) -| (s5.north);
\end{tikzpicture}
\end{bkfigure}
\vspace{1mm}
例題は「場面→何を無視するか→図と記号→使う式を選ぶ理由→式を立てる→解く→答えを確かめる」の7つの段で解きます。こまったら、例題にもどろう。

\bksubsection{2}{章のはじめの「なぜ？」と「道しるべ」}
\leavevmode 章の帯には、その章で答える「なぜ？」があります。帯のすぐ下には、こんな行があります。
\par\vspace{1.5mm}
\noindent\begin{minipage}{\linewidth}\bkmichishirube{5-12,8-8,8-9,電-1,電-2,熱-1}{電-5,電-9,電-16}\end{minipage}
数字の札（8-8）は数学シリーズ、漢字の札（電-1・熱-1）は物理シリーズの章です。わからない式は、数学の本や第1〜3巻（『力学』『熱学』『波動』）にもどって確かめよう。

\clearpage
\bksubsection{3}{式が出てきたら、3つの見方で読む}
\begin{phyomi}[$P=VI$]
電力 $P$（1秒に使うエネルギー）は、電圧 $V$ と電流 $I$ の積。\par\noindent\phdim{W ＝ V × A}\enspace{\small 左と右の単位がそろう（V＝J/C、A＝C/s）}
\end{phyomi}
\vspace{-2mm}
\begin{phmijika}
\qty{100}{V}で\qty{1200}{W}のドライヤーには\qty{12}{A}が流れ、1分で\qty{72}{kJ}の熱を出す。
\end{phmijika}
\vspace{-2mm}
\begin{phbaisuru}{電圧 $V$}
同じ電熱線なら電流も2倍になり、電力は4倍（$P=V^2/R$）。
\end{phbaisuru}

\bksubsection{4}{物理の本で足した囲み}
\noindent\ph@newbox{phmodel}{phModel}{モデル化の札}{無視したこと・仮定・くずれるとき}\hfill
\ph@newbox{phketa}{phKeta}{桁の感覚}{本物の値と、答えの桁の確かめ方}\par\vspace{1mm}
\noindent\ph@newbox{phstation}{phInk}{工学の駅}{この先の工学の科目と専門書のことば}\hfill
\ph@newbox{phflow}{bkTheme}{一般の式から、よく使う式へ}{条件をおいて、よく使う式を導く}\par\vspace{1mm}
\noindent\ph@newbox{phmada}{phMada}{まだわかっていないこと}{今の科学でも答えが出ていないこと}\hfill
\ph@newbox{phsaki}{phSaki}{この先の学び}{この本の範囲をこえる理論の入口}\par

\bksubsection{5}{★と3つの段}
★は内容の高さ（★中学理科、★★高校物理基礎、★★★高校物理、★★★★大学の初め、★★★★★大学の専門）。もくじの\phstarmeter{2}{4}は★2〜★4（ぬった★が章のはじめ、白ぬきの☆までが最も高い所）。3段の章は\textbf{目でわかる}→\textbf{手でわかる}→\textbf{式でわかる}で、「ここまでで十分」の旗まで読めばゴールです。

\bksubsection{6}{この本の4つのやくそく}
\begin{enumerate}[label={\bkcircled[bkNavy]{\arabic*}}, leftmargin=2.4\zw, itemsep=2pt, topsep=2pt]
\item 式を立てる前に、図をかいて、何を無視するかを書こう。
\item 答えに単位をつけ（Hz と rad/s、W と VA は別）、本物の値と桁を確かめよう。
\item 繰り返しの計算は、表計算やプログラムにまかせよう（シミュレーション）。
\item 実験は乾電池と磁石まで。家のコンセントや高い電圧にはさわらない。
\end{enumerate}
\clearpage\global\bk@decotrue}
\makeatother
